% Teaching manuscript. Edit this TeX source; Markdown is compatibility navigation only.

\StudyChapter{CS5489-Tutorial05}{Tutorial 5 · Face Detection with MLPs and Data Augmentation}
\StudyTarget{CS5489-Tutorial05__root-section}

\StudyRef{CS5489-Lecture05}{root}{Lecture 5：多层感知机与优化} · \StudyRef{CS5489-Lecture06}{root}{Lecture 6：CNN 与数据增强} · \StudyRepoRef{CS5489/course-notes/README.md}{}{课程目录} · \StudyRef{CS5489-Tutorial04}{root}{上一份：Tutorial 4} · \StudyRef{CS5489-Tutorial06}{root}{下一份：Tutorial 6}


\StudyHeading{1}{导读｜把 19×19 灰度图块展平成向量，为什么纯全连接网络很快遇到泛化天花板？}
\StudyTarget{CS5489-Tutorial05__overview}

Tutorial 5（对应 Canvas 原文件 \texttt{Tutorial5.ipynb}，共 55 个单元，配套数据包 \texttt{faces.zip} 与实景测试图 \texttt{nasa-small.png}）是全课程第一次使用 \textbf{PyTorch} 训练神经网络处理真实图像分类与滑动窗口目标检测任务。它的核心教学线索非常清晰：先将 $19\times 19\times 1$ 的人脸/非人脸小图块直接拉直成 $361$ 维向量喂给逻辑回归与多层感知机（MLP），亲眼看到“参数量庞大而训练图块仅约七千张”引发的严重过拟合；随后引入在线\textbf{数据增强（data augmentation）}观察测试准确率的跃升，最后用步长为 $4$ 的滑动窗口在宇航员合照上做整图人脸检测，为下一节\StudyRef{CS5489-Tutorial06}{root}{Tutorial 6 改用卷积神经网络（CNN）}埋下直观伏笔。

\begingroup\small\setlength{\tabcolsep}{3pt}\renewcommand{\arraystretch}{1.22}
\begin{longtable}{@{}>{\raggedright\arraybackslash}p{\dimexpr0.32\linewidth-1.50000\tabcolsep\relax}>{\raggedright\arraybackslash}p{\dimexpr0.16\linewidth-1.50000\tabcolsep\relax}>{\raggedright\arraybackslash}p{\dimexpr0.22\linewidth-1.50000\tabcolsep\relax}>{\raggedright\arraybackslash}p{\dimexpr0.30\linewidth-1.50000\tabcolsep\relax}@{}}
\toprule
原 Notebook 单元段 & 优先级 & 常见易错点 & 掌握要求 \\
\midrule\endfirsthead
原 Notebook 单元段 & 优先级 & 常见易错点 & 掌握要求 \\
\midrule\endhead
\bottomrule\endfoot
第 1–18 单元：读取 \texttt{faces.zip} 与划分验证集 & 核心必会 & 测试集子采样平衡、\texttt{Flatten} 形状 & 说清 $(6977,19,19,1)$ 拆分 $90\%/10\%$ 训练/验证的作用 \\
第 19–29 单元：PyTorch 逻辑回归与 MLP 训练 & 核心必会 & 参数量暴涨导致训练/验证曲线劈叉 & 解释为何宽深 MLP 在小样本扁平像素上易过拟合 \\
第 30–44 单元：像素噪声与几何数据增强 & 核心必会 & 标签保持不变性（禁 $180^\circ$ 倒置） & 区分高斯噪声、像素置零、对比度缩放、微旋与水平翻转的效果 \\
第 45–55 单元：\texttt{nasa-small.png} 滑窗检测 & 常规掌握 & 窗口步长、尺度固定与背景误报 & 解释固定 $19\times 19$ 滑窗为何只检同尺度人脸且产生密集假阳性 \\
\end{longtable}\endgroup


\StudyHeading{1}{1. 数据结构、训练/验证划分与 PyTorch 线性基线（第 1–24 单元）}
\StudyTarget{CS5489-Tutorial05__data-baseline}

原 Notebook 直接从不解压的 \StudyInline{faces.zip} 读取灰度图块：训练集共 $N_{\mathrm{train}} = 6{,}977$ 张，测试集对正负类做对称截断（每类最多 $472$ 张，共 $N_{\mathrm{test}} = 944$ 张，保证测试基准类别均衡 $50\%/50\%$）。每张图块尺寸为 $19\times 19\times 1$，拉直后输入维度为 $D = 19\times 19 = 361$。为调参选模型而不偷看测试集，第 16 单元用 \StudyInline{train\_\allowbreak{}test\_\allowbreak{}split(...,\allowbreak{} train\_\allowbreak{}size=0.9,\allowbreak{} test\_\allowbreak{}size=0.1,\allowbreak{} random\_\allowbreak{}state=4488)} 固定切出 $6{,}279$ 张训练样本（\StudyInline{vtrainI}）与 $698$ 张验证样本（\StudyInline{validI}）。

基线模型用 \texttt{nn.Sequential(nn.Flatten(), nn.Linear(361, 2))} 配合交叉熵损失训练二分类器。注意 PyTorch 实现细节：\texttt{nn.CrossEntropyLoss()} 内部已经整合了 \texttt{LogSoftmax} 与负对数似然，因此网络最后一层只需输出未归一化的两类原始对数几率（logits），切勿在送入 \texttt{CrossEntropyLoss} 前重复套一层 \texttt{Softmax}。


\StudyHeading{1}{2. 多层感知机（MLP）架构探索与过拟合诊断（第 25–29 单元）}
\StudyTarget{CS5489-Tutorial05__mlp}

将单层线性分类器扩展为带隐藏层与非线性激活（如 \texttt{ReLU}）的 MLP（例如 $361 \to 128 \to 64 \to 2$）后，首层权重矩阵立刻产生 $361\times 128 + 128 = 46{,}336$ 个可学习参数——可调参数量达到训练样本数（$6{,}279$）的七倍以上。同时，直接把 $19\times 19$ 图像拍扁成无序向量丢掉了二维空间邻域结构：同一只眼睛稍微向左平移 $2$ 个像素，落在完全不同的输入神经元下标上，MLP 必须重新从头记忆每一种平移位置。

因此训练曲线呈现出\StudyRef{CS5489-Lecture05}{monitoring}{Lecture 5 典型过拟合特征}：训练损失很快逼近 $0$、训练准确率接近 $100\%$，而验证损失在十余个 epoch 后见底回升。调优时需结合三项手段：（1）适度控制隐藏层宽度；（2）加入 \textbf{Dropout} 或 \textbf{L2 权重衰减（\texttt{weight\_decay}）}；（3）监控验证集准确率执行\textbf{早停（early stopping）}。


\StudyHeading{1}{3. 在线数据增强：给小数据集注入几何与光照不变性（第 30–44 单元）}
\StudyTarget{CS5489-Tutorial05__augmentation}

既然 $6{,}279$ 张图块不够喂饱 MLP，第 30–44 单元引入两类在每个训练 batch 动态采样的\textbf{数据增强（data augmentation）}（对应\StudyRef{CS5489-Lecture06}{early-stopping}{Lecture 6 正则化与增强专题}）：
\begin{itemize}
\item \textbf{像素强度扰动}：（1）叠加零均值高斯噪声 $\mathbf{x} + \mathcal{N}(0, \sigma^2)$（模拟传感器热噪声）；（2）随机像素置零腐蚀（pixel corruption / cutout，逼迫网络不依赖单一像素点）；（3）随机对比度缩放与亮度偏移 $a\mathbf{x} + b$（模拟明暗光照变化）。
\item \textbf{二维几何变换}：（1）\textbf{随机水平翻转（\texttt{RandomHorizontalFlip}）}——正脸左右高度对称，镜像翻转后百分之百仍是真实自然的正脸；（2）\textbf{小角度随机旋转（\texttt{RandomRotation(10)}，$\pm 10^\circ$）}与微小平移裁剪——模拟拍摄时头部轻微歪斜与框选抖动。
\end{itemize}

\textbf{原课核心思考题辨析（第 44 单元）：哪种增强提升最大？为什么？}实验表明，\textbf{水平翻转配合小角度旋转/平移及适度亮度对比度扰动}带来的测试集增益远大于单纯加高斯白噪声。原因在于：测试集与真实照片中的困难样本并不是“原图撒了一层雪花点”，而是人脸姿态略微偏转、框中心偏移一两个像素或左右打光不均；几何与光照增强精准填补了扁平 MLP 最缺乏的\textbf{局部平移/微旋与光照不变性先验}。反之，若把旋转角度设成 $\pm 180^\circ$ 或启用上下垂直翻转，制造出的倒立人脸与正立测试集分布背道而驰，反而会损害正立检测精度。此外启用强在线增强后，每次遍历看到的都是新变体，过拟合显著延后，需相应关闭过早触发的早停并适当增加训练轮数。


\StudyHeading{1}{4. 实景图像滑动窗口检测与局限分析（第 45–55 单元）}
\StudyTarget{CS5489-Tutorial05__sliding-window}

最后第 45–55 单元读取灰度合照 \StudyInline{nasa-\allowbreak{}small.png}，调用 \StudyInline{skimage.\allowbreak{}util.\allowbreak{}view\_\allowbreak{}as\_\allowbreak{}windows(\allowbreak{}testimg,\allowbreak{} (19, 19),\allowbreak{} step=4)} 以 $4$ 像素步长密集切出成千上万个 $19\times 19$ 候选窗口批量送入训练好的 MLP 预测，再将二值预测网格放大回原图尺寸叠加显示。观察可视化结果可总结两条直接通向后续课程的关键认识：
\begin{enumerate}
\item \textbf{高假阳性率（False Positives）}：分类器测试准确率哪怕达到 $85\%$，在整张图数千个背景窗口上扫一遍，仍会把大量衣领纹理、徽章或明暗交界误判为人脸（如 $15\%$ 误报率乘上 $2{,}000$ 个背景块就是约 $300$ 个误报亮点），需要更高的特异度、置信度阈值过滤以及非极大值抑制（NMS）。
\item \textbf{单尺度限制}：固定 $19\times 19$ 窗口只能检出恰好约 $19\times 19$ 像素大小的人脸；若照片里有近处大脸或远处小脸，传统滑窗必须构建多尺度图像金字塔（image pyramid），而在\StudyRef{CS5489-Lecture12}{detectors}{Lecture 12 目标检测（Faster R-CNN / YOLO）}中则会由特征金字塔与锚框端到端解决。
\end{enumerate}

\textbf{双语自测 / Check（对照原题问点）：}为什么在启用随机数据增强训练时，原 Notebook 提示先关闭默认早停（disable early stopping）？ / Why do we disable aggressive early stopping when training with random online data augmentation?

\begin{StudySolution}{答案 / Answer}
在线增强每轮随机生成新的旋转、翻转与噪声变体，训练初期损失下降更慢且验证曲线存在短期波动；若沿用未增强时很小的 \texttt{patience} 早停，网络还没充分见过丰富的增强样本就会过早终止，导致欠拟合。 / Online augmentation continuously synthesizes new transformed views, slowing initial convergence and adding small validation fluctuations; stopping too early prevents the network from learning across sufficient augmented variations.
\end{StudySolution}

来源：2026 学期 \texttt{Lecture5/Tutorial/Tutorial5.ipynb}（第 1–55 单元）。
